
\section{Emergent Behaviors within a Multi-agent Group}
\label{ssec:emerging_behaviors}

\begin{figure}[!ht]
    \centering
    \vspace{-1em}
    \includegraphics[width=\textwidth]{figs/files/minecraft-pipeline__.pdf}
    \vspace{-1.5em}
    \caption{An illustration of the collaborative process involving three agents crafting a bookshelf. The process begins with the decision-making and breaking down the goal into several sub-tasks, with each agent receiving an assignment. The execution results and the current environmental state are then passed to the evaluator. This process repeats until the goal of crafting a bookshelf is achieved.}
    \label{fig:minecraft-pipeline}
    \vspace{-0.5em}
\end{figure}
\begin{figure}[ht]
    \centering
    \includegraphics[width=\textwidth]{figs/files/minecraft-example.pdf}
    \vspace{-1.5em}
    \caption{Examples of the properties emerge in the agent interactions in Minecraft.}
    \label{fig:minecraft-example}
    \vspace{-1em}
\end{figure}
In the preceding section, the efficacy of \framework has been illustrated across a spectrum of tasks that necessitate multi-agent decision-making, especially for GPT-4-based agents. Our endeavor, however, surpasses just improvements on benchmark datasets. We delve deeper into emergent collaborative behaviors exhibited by agents within realistic, embodied AI contexts. Minecraft, a sandbox game, serves as an ideal platform for such exploration due to its intricate parallelisms with real-world dynamics. In the game, agents must not just execute tasks but also plan, coordinate, and adjust to evolving situations. We task agents with collaboratively crafting a variety of items, spanning from paper and paintings to books and bookshelves. A succinct figure showcasing three agents adeptly crafting a bookshelf can be viewed in~\cref{fig:minecraft-pipeline}. An elaborate visualization is placed at~\cref{sec:appendix-examples}, and details of the setups can be found in~\cref{sec:appendix-minecraft-details}.

By examining the decision-making process, we identify several emergent behaviors and categorize them into three aspects: \textit{volunteer}, \textit{conformity}, and \textit{destructive} behaviors. Note that these behaviors not necessarily only appear in Minecraft but also in previous experiments such as tool utilization.



\subsection{Volunteer Behaviors}
\label{ssec:emergent-volunteer}
Volunteer behaviors refer to actions intended to enhance the benefits of others in human society~\citep{omoto1995sustained,mowen2005volunteer}. We observe similar behaviors emerging in a multi-agent group as follows:  


\textbf{Time Contribution.} 
The agents are willing to contribute their unallocated time to enhance collaboration efficiency. As shown in the examples in~\cref{fig:minecraft-example} (1a), Alice and Bob need to collaboratively craft 2 paper, which necessitates three sugar canes as the raw material. Initially, Alice proposes that she will collect the sugar canes while Bob waits until the materials are ready. However, this plan is suboptimal, as it offers Bob spare time. Recognizing inefficiency, Bob suggests that both gather sugar canes concurrently, leading to expedited task completion. 

\textbf{Resource Contribution.}
Our analysis reveals that the agents are willing to contribute the possessed materials. As illustrated in \cref{fig:minecraft-example} (1b), at the end of the task crafting 2 paper, Alice has collected all the raw materials (sugar canes), whereas Bob possesses the crafting table essential for the paper's creation. In the decision-making stage, Alice suggests transferring her materials to Bob by dropping them on the ground. This enables Bob to utilize them for the intended crafting process. 



\textbf{Assistance Contribution.}
In the process of accomplishing tasks, we observe that agents, upon completing their individual assignments, actively extend support to their peers, thereby expediting the overall task resolution. As shown in \cref{fig:minecraft-example} (1c), Alice and Bob have successfully completed their assigned sub-tasks, while Charlie is still struggling to gather three leathers. During the collaborative decision-making phase, Alice and Bob propose to assist Charlie in gathering. 

These behaviors highlight how agents willingly contribute their capabilities and efforts to assist other agents, culminating in an accelerated achievement of their mutual goal.


\subsection{Conformity Behavior}
In human society, individuals tend to adjust their behavior to align with the norms or goals of a group~\citep{cialdini2004social, cialdini1998social}, which we refer to as \textit{conformity behavior}. We also observe similar behaviors within multi-agent groups. As shown in \cref{fig:minecraft-example} (2), all agents are asked to gather three pieces of leather. However, Charlie gets sidetracked and begins crafting items that do not contribute directly to the task. In the subsequent decision-making stage, Alice and Bob critique Charlie's actions. Charlie acknowledges his mistake and re-focuses on the mutual tasks. The conformity behavior enables agents to align with mutual goals as work progresses.










\subsection{Destructive behavior}
Additionally, we have also observed that agents may exhibit behaviors aimed at achieving greater efficiency, which could raise safety concerns. As depicted in \cref{fig:minecraft-example} (3a) and \cref{fig:minecraft-example} (3b), an agent occasionally bypasses the procedure of gathering raw materials and resorts to harming other agents or destroying an entire village library to acquire the necessary materials.

With advancements in autonomous agents, deploying them in real-world scenarios has become increasingly plausible. However, the emergence of hazardous behaviors could pose risks, especially when humans are involved in collaborative processes. Thus, designing strategies to prevent agents from adopting such hazardous behaviors is a critical area for future research.

